% =========================================================================
\section{Circuito 1: amplificador inversor}\label{sec:pl-inversor}
% =========================================================================
\begin{prelab}[Amplificador inversor]\label{pl:inversor}
Usando el modelo de amplificador operacional ideal visto en teoría, encuéntrese la ganancia en
tensión $A_V=\Vout/\Vin$ en función de las resistencias del circuito (\cref{fig:inversor}). Si
imponemos $R_1=1\,\mathrm{k}\Omega$, determínese el valor de $R_2$ para conseguir $A_V=-10$ a $1\,\mathrm{kHz}$.
\end{prelab}

% Figura: amplificador inversor (circuito 1) con tensiones y corrientes del análisis.
\begin{figure}[ht]
\centering
\begin{circuitikz}[circuito]
  \draw (0,0) node[op amp, ao] (A) {};
  % nudo inversor y red de realimentación
  \draw (A.-) -- ++(-0.6,0) coordinate (N);
  \draw (N) to[R, l_=$R_1$, i<^=$i_1$] ++(-2.4,0) coordinate (E);
  \draw (N) |- ++(0,1.5) coordinate (T);
  \draw (A.out) -- ++(0.7,0) coordinate (O);
  \draw (T) to[R, l=$R_2$, i>^=$i_2$] (T -| O) -- (O);
  % generador
  \draw (E) to[short, -o] ++(-0.3,0) coordinate (Ein) node[above, tens] {$\vin$};
  \draw (Ein) -- ++(-0.9,0) coordinate (G) to[sV, l_=$v_g$] ++(0,-2.0) node[ground] {};
  % entrada no inversora a tierra
  \draw (A.+) -- ++(-0.6,0) coordinate (P) -- ++(0,-0.45) node[ground] {};
  % carga y salida
  \draw (O) to[short, *-o] ++(1.0,0) coordinate (S) node[right, tens] {$\vout$};
  \draw (O) ++(0.5,0) to[R, l=$R_L$] ++(0,-2.0) node[ground] {};
  \fill (N) circle (1.6pt);
  % anotaciones del análisis
  \node[tens, anchor=east, font=\scriptsize] at ($(N)+(-0.08,0.75)$) {$\vm=0$};
  \draw[azulprof, line width=0.3pt] ($(N)+(-0.08,0.72)$) -- ($(N)+(-0.02,0.06)$);
  \node[tens, anchor=east, font=\scriptsize] at ($(P)+(-0.1,-0.25)$) {$\vp=0$};
\end{circuitikz}
\caption{Amplificador inversor (figura~2 del guion~\cite{GUION-S1}), con $R_1=1\,\mathrm{k}\Omega$,
$R_L=10\,\mathrm{k}\Omega$ y la alimentación omitida. Se indican las magnitudes del análisis: el nudo
inversor es una tierra virtual ($\vm=\vp=0$) y no absorbe corriente, de modo que $i_1=i_2$.}
\label{fig:inversor}
\end{figure}


\begin{resolucion}
Con los sentidos de corriente de la \cref{fig:inversor}:

\textbf{1. Tierra virtual.} La entrada $+$ está a tierra; por (R2),
\begin{equation}\label{eq:inv-vm}
  \vm=\vp=0 .
\end{equation}

\textbf{2. Ley de Ohm en $R_1$.}
\begin{equation}\label{eq:inv-i1}
  i_1=\frac{\vin-\vm}{R_1}=\frac{\vin}{R_1} .
\end{equation}

\textbf{3. Kirchhoff en el nudo inversor}, con $i_-=0$ por (R1):
\begin{equation}\label{eq:inv-kcl}
  i_2=i_1-i_-=\frac{\vin}{R_1} .
\end{equation}

\textbf{4. Ley de Ohm en $R_2$} ($i_2$ va del nudo inversor a la salida):
\begin{equation}\label{eq:inv-vout}
  \vout=\vm-i_2R_2=-\frac{R_2}{R_1}\,\vin .
\end{equation}

\textbf{Resultado.}
\begin{equation}\label{eq:inv-Av}
  \resultado{A_V=\frac{\vout}{\vin}=-\frac{R_2}{R_1}}
\end{equation}

\textbf{Diseño.} Con $A_V=-10$ y $R_1=1\,\mathrm{k}\Omega$ en \cref{eq:inv-Av}:
\begin{equation}\label{eq:inv-R2}
  R_2=\abs{A_V}\,R_1\quad\Longrightarrow\quad\resultado{R_2=10\,\mathrm{k}\Omega}
\end{equation}

\textbf{Comentario.} El signo indica contrafase ($180^{\circ}$). En el modelo ideal $A_V$ no depende de la
frecuencia (vale igual a $1\,\mathrm{kHz}$), ni de $A$ ni de $R_L$: solo del cociente de resistencias
(\cref{fig:inv-graficas}).
\end{resolucion}

% Figura: previsiones del amplificador inversor (A_V=-10, Vsat=9 V, f=1 kHz).
\begin{figure}[ht]
\centering
\begin{tikzpicture}
  % (a) zona lineal
  \begin{axis}[tema, name=a, width=0.38\textwidth, height=4.6cm,
      xmin=0, xmax=2, ymin=-10.5, ymax=10.5, ytick={-9,-5,0,5,9},
      xlabel={$t$ (ms)}, ylabel={Tensión (V)}, title style={font=\small},
      title={(a) $V_p=0{,}5\,\mathrm{V}$}]
    \addplot[asintota, domain=0:2, samples=2] {9};
    \addplot[asintota, domain=0:2, samples=2] {-9};
    \addplot[senal entrada, domain=0:2, samples=150] {0.5*sin(360*x)};
    \addplot[senal salida, domain=0:2, samples=150] {-5*sin(360*x)};
  \end{axis}
  % (b) característica de transferencia
  \begin{axis}[temacentro, name=b, at={(a.east)}, anchor=west, xshift=0.9cm,
      width=0.33\textwidth, height=4.6cm,
      xmin=-1.6, xmax=1.6, ymin=-11, ymax=11, xtick={-0.9,0.9}, ytick={-9,9},
      xticklabels={$-0{,}9$,$0{,}9$}, yticklabels={$-9$,$9$},
      xlabel={$\vin$ (V)}, 
      xlabel style={at={(axis description cs:1,0.5)}, anchor=south east, yshift=2pt, font=\scriptsize}, clip=false,
      title style={font=\small, yshift=10pt}, title={(b) $\vout(\vin)$}]
    \addplot[senal salida, domain=-1.6:1.6, samples=200] {max(-9,min(9,-10*x))};
    \addplot[punteado] coordinates {(0.9,0) (0.9,-9)};
    \addplot[punteado] coordinates {(-0.9,0) (-0.9,9)};
    \node[font=\scriptsize, anchor=west, text=rojoteo] at (axis cs:0.1,6.6) {pendiente $-10$};
    \node[font=\scriptsize, anchor=west] at (axis description cs:0.52,0.97) {$\vout$ (V)};
  \end{axis}
  % (c) recorte
  \begin{axis}[tema, at={(b.east)}, anchor=west, xshift=0.7cm, width=0.38\textwidth, height=4.6cm,
      xmin=0, xmax=2, ymin=-16, ymax=16, ytick={-15,-9,0,9,15}, yticklabels={},
      xlabel={$t$ (ms)}, title style={font=\small},
      title={(c) $V_p=1{,}5\,\mathrm{V}$}]
    \addplot[gray!60, densely dashed, domain=0:2, samples=150] {-15*sin(360*x)};
    \addplot[asintota, domain=0:2, samples=2] {9};
    \addplot[asintota, domain=0:2, samples=2] {-9};
    \addplot[senal entrada, domain=0:2, samples=150] {1.5*sin(360*x)};
    \addplot[senal salida, domain=0:2, samples=300] {max(-9,min(9,-15*sin(360*x)))};
    \node[font=\scriptsize, anchor=south east] at (axis cs:2,9) {$+\Vsat$};
  \end{axis}
\end{tikzpicture}
\caption{Previsiones para el amplificador inversor ($A_V=-10$, $\Vsat\approx9\,\mathrm{V}$,
$f=1\,\mathrm{kHz}$); en azul la entrada y en ámbar la salida. (a) En zona lineal la salida es la
entrada multiplicada por $-10$ (contrafase). (b) Característica de transferencia: recta de pendiente
$-10$ hasta $\abs{\vin}=\Vsat/10=0{,}9\,\mathrm{V}$ y saturación fuera. (c) Con $V_p=1{,}5\,\mathrm{V}$ la salida
lineal, en gris, superaría $\Vsat$ y la salida real queda recortada en $\pm\Vsat$ (\emph{clipping}).}
\label{fig:inv-graficas}
\end{figure}

